\documentclass[12pt]{article}
\usepackage[T2A]{fontenc}
\usepackage[utf8]{inputenc}
\usepackage[russian]{babel}
\usepackage{latexsym,amssymb,amsthm}
\usepackage{amsmath}
\usepackage{wrapfig}
\RequirePackage[pdftex]{graphicx}
%%\usepackage{array}

\usepackage{epsf,epic,eepic}
% \input ris.tex
% \input cells.sty
% \addrisoption{\risleftfalse}

\DeclareRobustCommand{\No}{\ifmmode{\nfss@text{\textnumero}}\else\textnumero\fi}

\oddsidemargin=0pt
\textwidth=159mm
\headheight=0pt
\headsep=0pt
\topmargin=0pt
\textheight=240mm

%%
%% Три точки для обозначения делимости
%%
\catcode`\@=11
\def\kratno{\mathrel{\smash{\lower.5ex\hbox{$\,\vdots\,$}}}}
% Это прямые три точки для обозначения делимости
\def\nekratno{\mathrel{\mathpalette\c@ncel\kratno}}
\def\c@ncel#1#2{\m@th\ooalign{$\hfil#1\mkern1mu/\hfil$\crcr$#1#2$}}
% Это перечеркнутые три точки для обозначения неделимости
\def\divis{\smash{\lower.1ex\hbox{\,\hbox{.}\kern-.22em\lower-.6ex\hbox{.}
\kern-.52em\lower-1.2ex\hbox{.}\,}} }
% А это наклонные...

\hyphenation{чис-ло чис-ла чис-лу чис-лом чис-ле чис-лам чис-ла-ми чис-лах
че-ты-рех-уголь-ник че-ты-рех-уголь-ни-ка пря-мо-уголь-ник тре-уголь-ни-ка
тре-уголь-ни-ков вне-впи-сан-ной вне-впи-сан-ных}

\def\signed(#1){{\unskip\nobreak\hfil\penalty50\hskip1em\hbox{}\nobreak\hfil
                 \rm(\small\it#1\/\rm)\parfillskip=0pt\finalhyphendemerits=0\par}}
        % macro for authorship's signature, from TeXbook
\def\be{\signed(С.~Берлов)}
\def\fp{\signed(Ф.~Петров)}
\def\frank{\signed(В.~Франк)}
\def\bah{\signed(Ф.~Бахарев)}
\def\guas{\signed(А.~Голованов)}
\def\as{\signed(А.~Смирнов)}
\def\dvk{\signed(Д.~Карпов)}
\def\dm{\signed(Д.~Максимов)}
\def\kk{\signed(К.~Кохась)}
\def\kuz{\signed(А.~Кузнецов)}
\def\hr{\signed(А.~Храбров)}
\def\ov{\signed(О.~Ванюшина)}
\def\io{\signed(О.~Иванова)}
\def\sop{\signed(Е.~Сопкина)}
\def\sviv{\signed(С.~Иванов)}
\def\suk{\signed(К.~Сухов)}
\def\dr{\signed(Д.~Ростовский)}
\def\pas{\signed(А.~Пастор)}
\def\azh{\signed(А.~Железняк)}
\def\vlas{\signed(Н.~Власова)}
\def\asol{\signed(А.~Солынин)}
\def\miv{\signed(М.~Иванов)}
\def\chuh{\signed(А.~Чухнов)}
\def\lev{\signed(Л.~Емельянов)}
\def\dsh{\signed(Д.~Ширяев)}
\def\zh{\signed(жюри)}
\def\bdzh{\signed(О.~Бадажкова)}

\sloppy
\pagestyle{empty}

\def\q#1.{{\bf #1. }}
\def\dotline{\smallskip\hbox to \hsize{\dotfill}\medskip}

\begin{document}

\centerline{\sc САНКТ-ПЕТЕРБУРГСКАЯ}
\centerline{\sc ОЛИМПИАДА ШКОЛЬНИКОВ 2021 года ПО МАТЕМАТИКЕ}
\centerline{\sc II тур. \ 6 класс.}
\medskip
\hrule
\bigskip

\q1. Саша задумал три натуральных числа. Он выписал на доску все общие делители первых двух чисел, затем выписал все общие делители второго и третьего, и наконец выписал все общие делители первого и~третьего. Всего оказалось выписано 18 чисел. Четыре из них~--- 3, 3, 4, 4. Найдите остальные 14 выписанных чисел. 
\kuz

\q2. В зале кинотеатра кресла стоят в несколько рядов (ряды могут различаться по количеству кресел).
Общая вместимость зала~--- 1000 зрителей. На сеанс пришло 515 зрителей, и 
их удалось рассадить так, что никакие два зрителя не оказались в~одном 
ряду на соседних креслах. Докажите, что в каком-то ряду не больше 33~кресел.
\asol

\q3. Шахматный конь обошел доску $8 \times 8$, побывав на каждой клетке по одному разу. Клетки пронумеровали вдоль маршрута коня числами от 1 до 64. После этого на клетку 1 пришла собака, и~стала искать коня по следующему правилу. Оказавшись на клетке с~номером $n$, она обнюхивает соседние по стороне клетки и находит среди них клетку с~наибольшим номером. Собака переходит на эту клетку, если ее номер больше $n$, и остается на месте, если номер меньше $n$. Известно, что собака нашла коня, то есть пришла на 64-ю клетку. Докажите, что она переходила с одной клетки на другую не более 19 раз.
\io

\q4. Дано натуральное число $n$. На доске выписаны все положительные 
несократимые дроби со знаменателем $n$, меньшие 10. Докажите, что их сумма 
не равна 123450. 

\dotline

\centerline{Олимпиада 2021 года. II тур. 6 класс. Выводная аудитория.}
\smallskip

\q5. Оля с Димой играют в игру. Сначала Оля выкладывает в ряд в некотором 
порядке $2n$ карточек с числами от 1 до $2n$ числами вверх. Затем 
они ходят по очереди, начиная с Димы. За один ход игрок забирает себе одну 
из крайних карт. Какое наибольшее количество подряд идущих чисел Оля может 
собрать у себя вне зависимости от действий Димы?  
\signed(О.~Бадажкова, Н.~Власова)

\q6. В большой таблице отмечено 2020 клеток (в каждом столбце и каждой строке отмечено не более одной клетки), при этом 
в~верхней строке и в левом столбце отмеченных клеток нет.
Сережа хочет покрасить эти клетки в красный и синий цвет так,
чтобы справа от любого столбца таблицы количество красных клеток отличалось от количества синих не более чем на 1,
и снизу от любой строки количество красных клеток отличалось от количества 
синих тоже не более чем на~1. 
Верно ли, что при любом расположении клеток Сережа сможет это сделать?


\clearpage

\centerline{\sc САНКТ-ПЕТЕРБУРГСКАЯ}
\centerline{\sc ОЛИМПИАДА ШКОЛЬНИКОВ 2021 года ПО МАТЕМАТИКЕ}
\centerline{\sc II тур. \ 7 класс.}
\medskip
\hrule
\medskip

\q1. Костя выписал в некотором порядке все дроби с числителем 2021 и 
знаменателями от 1 до 2020. Оказалось, что сумма первой и~второй дроби 
равна их произведению, сумма третьей и четвертой равна их произведению 
и т.\,д., сумма 2017-й и 2018-й дроби равна их произведению. Докажите, что
сумма 2019-й и 2020-й дроби тоже равна их произ\-ве\-дению.
\signed(К.~Кохась, Н.~Власова)

\q2. В треугольнике $ABC$ проведена медиана $AK$. На стороне $AC$ отмечена 
точка $N$ так,  что $\angle AKN = 90^\circ$. Докажите, что на плоскости найдется 
точка $L$ такая, что $AB = AL$, $BK=KL$ и $NL = NC$.
\kuz

\q3. Дана прямоугольная таблица $14 \times n$. В её левой верхней клетке стоит 
фишка. Саша и Оля по очереди (начинает Саша) передвигают фишку на одну 
клетку вправо или вниз. Будем говорить, что произошел поворот, если направление 
хода игрока отличалось от направления сделанного перед этим хода соперника. 
Если к моменту попадания фишки в~правый  нижний угол произошло чётное число 
поворотов, то побеждает Саша, иначе~--- Оля. При каких значениях $n$ Оля может 
обеспечить себе победу независимо от действий Саши?
\bdzh

\q4. В ряд выписано 300 натуральных чисел. Каждое число, начиная с~третьего,
равно произведению двух предыдущих. Сколько точных квадратов может быть 
среди этих чисел? (Приведите все ответы и~докажите, что других нет.)
\signed(Н.~Филонов)

\dotline

\centerline{Олимпиада 2021 года. II тур. 7 класс. Выводная аудитория.}
\smallskip

\q5. На плоскости даны 5000 точек, никакие две из которых не лежат на одной 
вертикальной или горизонтальной прямой. Сережа хочет так расставить в этих 
точках числа $1$ и $-1$, чтобы сумма всех чисел, расположенных слева от 
любой вертикальной прямой, была равна $0$, $1$ или $-1$, и сумма всех 
чисел, расположенных сверху от любой горизонтальной прямой, тоже была 
равна $0$, $1$ или $-1$. Верно ли, что при любом расположении точек он 
сможет этого добиться?
\be

\q6. Найдите наименьшее значение выражения $$\left[\frac{7(a+b)}{c}\right] + 
\left[\frac{7(a+c)}{b}\right]+ \left[\frac{7(b+c)}{a}\right],$$ где $a$, $b$ и $c$~--- 
произвольные натуральные числа. 
(Как обычно, через $[x]$ обозначается целая часть числа $x$,
т.\,е.\ наибольшее целое число, не превосходящее $x$.)
\hr

\q7. Внутри выпуклого четырехугольника $ABCD$ нашлись такие точки $E$ и $F$, что треугольники $ABE$ и $CDF$~--- равносторонние. Докажите, что сумма периметров треугольников $ABF$ и $CDE$ не меньше, чем периметр четырехугольника $ABCD$.
\kuz

\clearpage

\centerline{\sc САНКТ-ПЕТЕРБУРГСКАЯ}
\centerline{\sc ОЛИМПИАДА ШКОЛЬНИКОВ 2021 года ПО МАТЕМАТИКЕ}
\centerline{\sc II тур. \ 8 класс.}
\medskip
\hrule
\bigskip

\q1. Существует ли набор, состоящий из 2021 различного натурального числа, 
обладающий следующим свойством: если выбрать из него любое число $a$, 
остальные 2020 чисел можно разбить на пары так, чтобы $a$ делилось на 
разность чисел в каждой паре? 
%(КТ)

\q2. На сторону $AB$ остроугольного треугольника $ABC$ опущена высота $CH$, и на отрезке $AH$ отмечена точка $D$ так, что $AD = BH$. Докажите, что расстояние между серединами отрезков $BH$ и $CD$ вдвое меньше $AC$.
\kuz

\q3. Сизиф таскает на гору по одному камню в день. Первый и~второй камень весят по 1 кг, а вес каждого  следующего камня равен либо сумме, либо модулю разности весов двух предыдущих камней (каждый раз на выбор Сизифа, однако камней нулевого веса не бывает). Зевс назвал Сизифу два взаимно простых натуральных  числа $A$~и~$B$, и~разрешил ему уйти на заслуженный  отдых, если однажды тот отнесет на гору камень веса  $A$, а ровно через год (спустя 365 дней) после этого~--- камень веса $B$.  Докажите, что Сизиф сможет действовать так, чтобы когда-нибудь уйти на заслуженный отдых.  
\io

\q4. На берегу реки стоят 10 шейхов, каждый вместе с гаремом из 100 жён. Также у берега стоит $n$-местная яхта. По закону женщина не должна находиться на одном берегу, на яхте, или даже на пересадке с~мужчиной, если рядом нет её мужа. При каком наименьшем $n$ все шейхи с~женами смогут переправиться на другой берег, не нарушив закон?
\asol

\dotline

\centerline{Олимпиада 2021 года. II тур. 8 класс. Выводная аудитория.}
\smallskip

\q5. Найдите наименьшее значение выражения $$\left[\frac{8(b+c)}{a}\right] + \left[\frac{8(c+a)}{b}\right]+ \left[\frac{8(a+b)}{c}\right],$$
где $a$, $b$ и $c$~--- произвольные натуральные числа. (Как обычно, через $[x]$ обозначается целая часть числа 
$x$, т.\,е.\ наибольшее целое число, не превосходящее $x$.)
\hr

\q6. Назовем множество натуральных чисел, не превосходящих $10^{100}$, {\it консервативным}, если ни один из его элементов не является делителем произведения всех остальных элементов, и {\it прогрессивным}, если вместе с любым своим элементом оно содержит все кратные ему числа, не превосходящие $10^{100}$. Каких множеств больше: прогрессивных или консервативных?
\dsh

\q7. Дан выпуклый четырехугольник $ABCD$. Точка $A'$ такова, что $\angle ABA'=\angle ADA'=90^\circ$. Аналогично определяются точки $B', C'$ и $D'$. При этом $A'$ и $C'$ лежат внутри $ABCD$, а $B$ и $D$~--- внутри выпуклого четырехугольника $A'B'C'D'$. Докажите, что площади четырехугольников $ABCD$ и $A'B'C'D'$ равны.
\kuz

\end{document}
